% !TeX root = ../../Book4.tex
% 纲要标题：## 附录 G　Operad 与 Koszul 对偶初步
% * 状态：初稿已完成；新增范围尚未正式审核。

\chapter{Operad 与 Koszul 对偶初步}
\label{b4:app:G}
\BookChapterNumber{b4:app:G}

\begin{chapterabstract}
	乘法、括号及它们的复合可以先作为抽象运算研究，再作用于具体向量空间。Operad 用输入置换和代入组织这些运算；装饰根树构造自由对象，关系商给出结合、交换、Lie 与 Poisson 结构。二元二次对偶把关系换成带符号配对的零化空间，产生结合结构的自对偶及交换与 Lie 结构之间的对偶。树的 bar 复形则进一步记录高次相容条件，使二次对偶的构造与 Koszul 性的同调要求得以区分。
\end{chapterabstract}

\begin{learninggoals}
	\begin{itemize}
		\item 写出运算的输入重排和插入复合，区分一元恒等运算与代数的乘法单位。
		\item 用装饰根树构造自由 operad，并用泛性质给具体向量空间定义代数结构。
		\item 从生成元与关系识别 Ass、Com、Lie 与 Poisson 代数，核算三输入运算的维数。
		\item 计算二次关系的带符号零化空间，说明悬移为何改变输入交换与插入的符号。
		\item 写出树 bar 复形的边界及低次同调，区分二次对偶与需要额外证明的 Koszul 性。
	\end{itemize}
\end{learninggoals}

对一个双线性乘法 \(\mu\)，三次相乘有两种运算表达式
\[
	\mu(\mu(a,b),c),\qquad \mu(a,\mu(b,c)).
\]
它们在每个结合代数中相等，在没有结合律的代数中却可能不同。对四个输入，可以继续把这两种表达式代入更大的表达式；原来的一条关系就必须与所有这些代入相容。自由结合代数用字记录元素的乘积，而这里要先记录运算本身，以及运算之间的关系，然后才把它们作用于具体元素。

\ProseParagraphBreak
研究 Lie 括号时，基本关系改为反对称性与 Jacobi 恒等式；同时保留乘法和括号，还要加入 Leibniz 关系。树记录这些复合的嵌套方式。对关系空间作带符号的对偶，交换乘法的结合关系恰好变成括号的 Jacobi 关系；这不是两个公式的偶然相似，而是二次对偶中的一次线性代数计算。

\ProseParagraphBreak
本附录统一取特征零域，主要讨论普通向量空间及只含正元数的 operad；一元恒等运算保留，零元运算不列入。因此所编码的结合乘法不预先附带单位。所需基础是张量积、对称群作用、线性对偶及复形的同调；分次交换符号取\cref{b4:exm:graded-braiding}中的约定，Poisson 算例承接\subsecref{b4:subsec:1304:04}。二次对偶及经典例子可对照 Ginzburg--Kapranov 的正式论文\cite[Sections 1.2, 2.1]{GinzburgKapranov1994KoszulOperads}；最后的 Koszul 性结论给出证明概要，注明一般同调判据和标准消解比较所承担的步骤\cite[Sections 4.1--4.2]{GinzburgKapranov1994KoszulOperads}。

% !TeX root = ../../Book4.tex
% 纲要标题：### G.1 Operad 与代数结构
% * 状态：初稿已完成；新增范围尚未正式审核。
% #### G.1.1 对称 operad 与复合
% #### G.1.2 End operad 与 operad 上的代数
% #### G.1.3 自由构造与树形复合

\section{Operad 与代数结构}
\label{b4:sec:G01}

一个双线性乘法 \(\mu:V\otimes V\rightarrow V\) 不仅给出二元运算，还给出 \(\mu(\mu(a,b),c)\)、\(\mu(a,\mu(b,c))\) 等三元运算。结合律要求这两个复合相等；交换律比较同一运算在交换输入后的取值。若把乘法换成 Lie 括号，仍要讨论输入的置换和运算的代入，只是要求满足的关系变成反对称性与 Jacobi 恒等式。Operad 同时记录各种元数的运算、输入的重排和运算之间的代入，使这些代数结构可以用同一种语言描述。

\subsection{对称 operad 与复合}
\label{b4:subsec:G01:01}

以下在特征为零的域 \(k\) 上讨论普通向量空间。一个 \(n\) 元运算有 \(n\) 个输入和一个输出。这里允许一元恒等运算，只考虑 \(n\geq1\)，不放入零元运算。零元运算用于指定一个元素，例如乘法单位；一元恒等运算则把每个输入原样送出。这两种单位的作用不同。

\ProseParagraphBreak
用 \(1,\ldots,n\) 给输入编号，可以把它们的重排写成对称群 \(S_n\) 的作用。但在第 \(i\) 个输入处代入一个新运算以后，其余输入的编号会移动。为把等变条件写清楚，先允许输入由任意非空有限集 \(I\) 标记；不同集合之间的双射只改变标记。

\begin{definition}\label{b4:def:symmetric-sequence-operad}
	设 \(k\) 为域。对每个非空有限集 \(I\) 给定一个 \(k\) 向量空间 \(E(I)\)，对每个双射 \(\alpha:I\rightarrow J\) 给定线性同构
	\[
		\alpha^*:E(J)\longrightarrow E(I),
	\]
	满足 \(1_I^*=1_{E(I)}\) 以及 \((\beta\alpha)^*=\alpha^*\beta^*\)。这样的数据称为一个\term[duichen-xulie]{对称序列}[symmetric sequence]。记 \([n]=\{1,\ldots,n\}\)、\(E(n)=E([n])\)，并规定
	\[
		e\cdot\sigma=\sigma^*e\qquad(e\in E(n),\ \sigma\in S_n).
	\]
	置换乘法 \(\sigma\tau\) 表示先施行 \(\tau\)、再施行 \(\sigma\)，因而 \((e\cdot\sigma)\cdot\tau=e\cdot(\sigma\tau)\)，这是右 \(S_n\) 作用。
\end{definition}

对称序列也可只用一列右 \(k[S_n]\) 模 \(E(n)\) 给出。要恢复 \(E(I)\)，取所有 \(e\in E(n)\) 和编号双射 \(\ell:[n]\rightarrow I\) 的形式记号 \((e,\ell)\)，对 \(e\) 线性化，并加入关系
\[
	(e\cdot\sigma,\ell)=(e,\ell\sigma^{-1}).
\]
于是改变编号与改变右作用后的装饰表示同一个元素；双射 \(\alpha:I\rightarrow J\) 的拉回把 \((e,\ell)\) 送到 \((e,\alpha^{-1}\ell)\)。任选一个编号即可把这个空间识别为 \(E(n)\)，而上述关系保证识别之间正由规定的右作用转换。

\begin{definition}\label{b4:def:symmetric-operad}
	设 \(k\) 为域，\(P\) 是一个对称序列。对非空有限集 \(I,J\) 和 \(i\in I\)，给定双线性映射
	\[
		\circ_i:P(I)\otimes P(J)\longrightarrow P((I\setminus\{i\})\sqcup J),
		\qquad p\otimes q\longmapsto p\circ_i q,
	\]
	称为\term[charu-fuhe]{插入复合}[partial composition]。不交并中的输入均保留所属集合的标记；连续插入后按这些标记识别共同的输入集。另给定每个单元素输入集上的元素 \(\mathbf1\)，要求它们在重标号下互相对应，并满足以下条件。
	\begin{enumerate}
		\item 对 \(p\in P(I)\)、\(q\in P(J)\)、\(r\in P(K)\) 及 \(j\in J\)，有
		\[
			(p\circ_i q)\circ_j r=p\circ_i(q\circ_j r).
		\]
		若 \(i,i'\in I\) 且 \(i\ne i'\)，则
		\[
			(p\circ_i q)\circ_{i'}r=(p\circ_{i'}r)\circ_i q.
		\]
		\item 插入一元恒等运算不改变原运算：\(p\circ_i\mathbf1=p\)、\(\mathbf1\circ_1p=p\)，其中更换单元素输入标记后的空间按双射识别。
		\item 给定双射 \(\alpha:I\rightarrow I'\)、\(\beta:J\rightarrow J'\)，令
		\[
			\gamma:(I\setminus\{i\})\sqcup J\longrightarrow (I'\setminus\{\alpha(i)\})\sqcup J'
		\]
		在两个部分上分别等于 \(\alpha\)、\(\beta\)。对 \(p\in P(I')\)、\(q\in P(J')\)，要求
		\[
			(\alpha^*p)\circ_i(\beta^*q)=\gamma^*(p\circ_{\alpha(i)}q).
		\]
	\end{enumerate}
	这样的数据称为一个\term[operad]{对称 operad}[symmetric operad]，简称 operad。它的态射 \(P\rightarrow Q\) 是与输入重标号、插入复合及 \(\mathbf1\) 相容的一族线性映射 \(P(I)\rightarrow Q(I)\)。
\end{definition}
\symidx{\(P(n)\)：operad 的 \(n\) 元运算空间}
\symidx{\(p\circ_iq\)：operad 中在第 \(i\) 个输入处的插入复合}

第一条结合等式比较嵌套的代入，第二条比较在两个不同输入处的代入；等变条件则说明先重排输入再代入，与先代入再重排最终输入得到同一运算。若使用数字编号，约定 \(q\in P(m)\) 插在 \(p\in P(n)\) 的第 \(i\) 个位置后，它的输入占据 \(i,\ldots,i+m-1\)，原来在其后的输入顺次后移。两条结合等式便写成
\[
	\begin{aligned}
		(p\circ_iq)\circ_{i+j-1}r&=p\circ_i(q\circ_jr) &&(1\leq j\leq m),\\
		(p\circ_iq)\circ_{j+m-1}r&=(p\circ_jr)\circ_iq &&(1\leq i<j\leq n).
	\end{aligned}
\]
同时代入 \(q_1,\ldots,q_n\) 的结果记为 \(p(q_1,\ldots,q_n)\)：可以从右向左逐个插入，各次代入的输入次序仍按原来的区块排列。不同输入处的结合等式保证结果不依赖代入这些区块的先后次序。

\ProseParagraphBreak
一元空间 \(P(1)\) 在 \(\circ_1\) 下是含幺结合代数，但不要求它只有一维。若 \(k\rightarrow P(1)\)、\(c\mapsto c\mathbf1\) 为同构，则称 \(P\) 为\term[liantong-operad]{连通 operad}[connected operad]。后面的二元生成例子满足这一条件，所有向量空间的自同态运算则一般不满足。无零元运算的约定也意味着：仅从 \(\mathbf1\in P(1)\) 不能得到某个向量 \(1_V\in V\)。

\subsection{End operad 与 operad 上的代数}
\label{b4:subsec:G01:02}

要使抽象符号 \(p\in P(n)\) 真正作用在向量空间 \(V\) 上，必须把它送到一个多线性映射，并要求符号的代入变成映射的代入。全部多线性映射本身就组成一个 operad。

\begin{proposition}\label{b4:prop:endomorphism-operad}
	设 \(k\) 为域，\(V\) 为 \(k\) 向量空间。令
	\[
		\operatorname{End}_V(n)=\operatorname{Hom}_k(V^{\otimes n},V)\qquad(n\geq1).
	\]
	对 \(f\in\operatorname{End}_V(n)\)、\(g\in\operatorname{End}_V(m)\)，规定
	\[
		(f\cdot\sigma)(v_1,\ldots,v_n)=f(v_{\sigma^{-1}(1)},\ldots,v_{\sigma^{-1}(n)}),
	\]
	\[
		\begin{aligned}
			&(f\circ_i g)(v_1,\ldots,v_{n+m-1})\\
			&\qquad=f(v_1,\ldots,v_{i-1},g(v_i,\ldots,v_{i+m-1}),v_{i+m},\ldots,v_{n+m-1}).
		\end{aligned}
	\]
	以 \(1_V\) 为一元恒等运算，这些空间组成 \(V\) 的\term[zidongtai-operad]{自同态 operad}[endomorphism operad]，记为 \(\operatorname{End}_V\)。
\end{proposition}
\symidx{\(\operatorname{End}_V(n)\)：向量空间 \(V\) 上的 \(n\) 元多线性运算}
\begin{proof}
	由张量积的泛性质，多线性映射与 \(V^{\otimes n}\rightarrow V\) 的线性映射一一对应。上式的代入仍对每个输入线性，故给出所需映射。在右作用的复合中，第 \(i\) 个输入最后变为 \(v_{\tau^{-1}\sigma^{-1}(i)}\)，所以
	\[
		((f\cdot\sigma)\cdot\tau)(v_1,\ldots,v_n)
		=f(v_{(\sigma\tau)^{-1}(1)},\ldots,v_{(\sigma\tau)^{-1}(n)}).
	\]
	这正是 \(f\cdot(\sigma\tau)\)，而恒等置换不改变映射。

	对有限输入集 \(I\)，也可把 \(\operatorname{End}_V(I)\) 看成以 \(I\) 标记变量的多线性映射空间。双射 \(\alpha:I\rightarrow I'\) 的拉回把输入 \((v_i)_{i\in I}\) 交给原映射的第 \(j\) 个变量 \(v_{\alpha^{-1}(j)}\)。因此两种顺序的重标号与代入，都把相同的向量交给相同的变量，满足等变条件。

	嵌套的结合式两边都先把 \(r\) 的值送入 \(g\) 的指定输入，再把所得值送入 \(f\)；在不同输入处的结合式两边则分别把 \(g\) 和 \(r\) 的值送入 \(f\) 的那两个输入。对每组向量求值所得结果分别相等，故两种结合式都作为多线性映射相等。插入 \(1_V\) 原样返回其输入，给出单位律。
\end{proof}

特别地，\(\operatorname{End}_V(1)=\operatorname{End}_k(V)\) 包含 \(V\) 的全部线性自同态，而不只是标量倍的恒等映射。这个例子说明一般 operad 的一元运算空间不能预先限制为 \(k\)。

\begin{definition}\label{b4:def:operad-algebra}
	设 \(k\) 为域，\(P\) 为 \(k\) 上的 operad。一个 \(k\) 向量空间 \(V\) 连同 operad 态射
	\[
		\rho:P\longrightarrow\operatorname{End}_V
	\]
	称为一个 \(P\) 代数，也称为 \(P\) 上的\term[operad-daishu]{代数}[algebra over an operad]。记 \(p_V=\rho(p)\)。从 \(P\) 代数 \(V\) 到 \(W\) 的同态是线性映射 \(h:V\rightarrow W\)，满足
	\[
		h\bigl(p_V(v_1,\ldots,v_n)\bigr)=p_W(h(v_1),\ldots,h(v_n))
		\qquad(p\in P(n)).
	\]
\end{definition}

若 \(P\) 由一个二元运算 \(\mu\) 及关系 \(\mu\circ_1\mu=\mu\circ_2\mu\) 生成，那么 \(P\) 代数中的二元乘法满足结合律。这里的一元元素 \(\mathbf1\) 只要求 \(\rho(\mathbf1)=1_V\)，没有要求这个结合乘法存在单位。下一节将把这些生成元与关系写成具体 operad。

\begin{proposition}\label{b4:prop:free-operad-algebra}
	设 \(k\) 为域，\(P\) 为只含正元数的 \(k\) 线性 operad，\(V\) 为 \(k\) 向量空间。令 \(S_n\) 在 \(V^{\otimes n}\) 上左作用为
	\[
		\sigma\cdot(v_1\otimes\cdots\otimes v_n)
		=v_{\sigma^{-1}(1)}\otimes\cdots\otimes v_{\sigma^{-1}(n)}.
	\]
	则
	\[
		P(V)\coloneqq\bigoplus_{n\geq1}P(n)\otimes_{k[S_n]}V^{\otimes n}
	\]
	自然成为一个 \(P\) 代数。映射 \(j:V\rightarrow P(V)\)、\(j(v)=[\mathbf1\otimes v]\) 使它成为 \(V\) 上的\term[ziyou-operad-daishu]{自由 \(P\) 代数}[free algebra over an operad]：对任意 \(P\) 代数 \(A\) 和线性映射 \(f:V\rightarrow A\)，存在唯一的 \(P\) 代数同态 \(\widetilde f:P(V)\rightarrow A\) 满足 \(\widetilde f j=f\)。
\end{proposition}
\symidx{\(P(V)\)：由向量空间 \(V\) 生成的自由 \(P\) 代数}
\begin{proof}
	张量积中的平衡关系为
	\[
		[(p\cdot\sigma)\otimes(v_1\otimes\cdots\otimes v_n)]
		=[p\otimes(v_{\sigma^{-1}(1)}\otimes\cdots\otimes v_{\sigma^{-1}(n)})].
	\]
	两边正表示同一运算改变输入编号前后的写法。对 \(p\in P(r)\) 及各个代表元 \([q_j\otimes\mathbf v_j]\)，其中 \(q_j\in P(n_j)\)、\(\mathbf v_j\in V^{\otimes n_j}\)，定义
	\[
		p_{P(V)}([q_1\otimes\mathbf v_1],\ldots,[q_r\otimes\mathbf v_r])
		=[p(q_1,\ldots,q_r)\otimes(\mathbf v_1\otimes\cdots\otimes\mathbf v_r)].
	\]
	在某个 \(q_j\) 处使用平衡关系，相当于只重排最终输入的第 \(j\) 个区块。插入复合的等变条件把这一置换移到右侧的张量上，两边仍给出同一商类。对不同的直和分量作多线性延拓，便得到良定义的运算。\(p\) 的输入置换同时重排这些区块；同一个等变条件给出 \(P\) 代数所需的右作用相容性。嵌套运算时，无论分两层还是一次代入，各 \(\mathbf v_j\) 的最终排列相同，运算部分由 operad 的结合律相等；单位条件也由 \(\mathbf1\) 的单位律得到。因此这些运算定义了 \(P\) 代数结构。

	规定
	\[
		\widetilde f([p\otimes(v_1\otimes\cdots\otimes v_n)])
		=p_A(f(v_1),\ldots,f(v_n)).
	\]
	\(A\) 的输入等变性保证它满足平衡关系，故这是线性映射；\(A\) 的代入相容性又说明它保持刚刚定义的全部 \(P\) 运算。由于结构映射满足 \(\rho(\mathbf1)=\operatorname{id}_A\)，有 \(\widetilde f j=f\)。相应的唯一分解可写成
	\[
		\begin{tikzcd}[column sep=large]
			V\arrow[r,"j"]\arrow[dr,"f"']&P(V)\arrow[d,dashed,"\exists!\,\widetilde f"]\\
			&A.
		\end{tikzcd}
	\]
	每个代表元都可写为
	\[
		[p\otimes(v_1\otimes\cdots\otimes v_n)]
		=p_{P(V)}(j(v_1),\ldots,j(v_n)).
	\]
	所以一个保持 \(P\) 运算且在 \(j(V)\) 上等于 \(f\) 的映射必须具有上面的公式，唯一性随之成立。
\end{proof}

\subsection{自由构造与树形复合}
\label{b4:subsec:G01:03}

自由 \(P\) 代数把已经规定的运算施加在生成向量上。另一种自由构造还没有预先规定运算之间的关系：给定若干基本运算，先把它们所有可能的代入保留下来，再按需要加入结合律或 Jacobi 恒等式。树正好记录一个复合运算由哪些基本运算、按怎样的嵌套次序组成。

\ProseParagraphBreak
本节的根树有一条输出边，每个顶点的输出朝向树根；进入顶点的边至少一条，没有顶点输出的输入边称为叶边。每条叶边用输入集 \(I\) 的一个元素标记，且全部标记恰好使用一次。也允许只有一条边、没有顶点的树，它有一个输入并将其直接连到输出。若顶点 \(v\) 的输入边集为 \(\operatorname{In}(v)\)，就用 \(E(\operatorname{In}(v))\) 中的一个元素装饰该顶点。这里 \(E\) 只是对称序列，还没有复合。

\begin{definition}\label{b4:def:tree-free-operad}
	设 \(k\) 为域，\(E\) 为 \(k\) 向量空间中的对称序列。令 \(\mathcal F(E)(I)\) 为上述有限装饰根树张成的空间，对每个顶点的装饰施加线性关系，并把保持根、叶标记及装饰的树同构识别。等价地，每棵树的装饰空间是
	\[
		\bigotimes_{v\in\operatorname{Vert}(T)}E(\operatorname{In}(v)),
	\]
	再对树的同构作识别；没有顶点的树的装饰空间约定为 \(k\)。这称为由 \(E\) 生成的\term[ziyou-operad]{自由 operad}[free operad]，其复合将一棵树的输出边接到另一棵树的指定叶边，接缝处不增加新顶点。
\end{definition}
\symidx{\(\mathcal F(E)\)：由对称序列 \(E\) 生成的自由 operad}

若把每个顶点的输入临时排成一列，则顶点装饰可写为 \(e\in E(m)\)。更换这一排列必须同时按右 \(S_m\) 作用更换装饰，正是前面的关系 \((e\cdot\sigma,\ell)=(e,\ell\sigma^{-1})\)。因此图中的左右位置只帮助读出一次选定的编号；在一般对称 operad 中，它并不是额外固定的平面结构。真正保留无对称二元运算的左右输入时，则可用正则右 \(S_2\) 模作生成元，把左右次序也记录在装饰中。

\begin{theorem}\label{b4:thm:free-operad}
	设 \(k\) 为域，\(E\) 为 \(k\) 向量空间中的对称序列。上述嫁接使 \(\mathcal F(E)\) 成为 operad。用单顶点树给出映射 \(\iota:E\rightarrow\mathcal F(E)\)。对任意 operad \(Q\) 和对称序列映射 \(\theta:E\rightarrow Q\)，存在唯一 operad 态射 \(\widehat\theta:\mathcal F(E)\rightarrow Q\) 满足 \(\widehat\theta\iota=\theta\)。
\end{theorem}
\begin{proof}
	输入重标号只改变叶标记。嫁接两棵树时，原顶点的输入边与接缝后的输入边有自然双射，故装饰可沿这些双射转移。树同构经过嫁接仍给出树同构，顶点装饰的线性关系也被保留，所以嫁接给出良定义的双线性映射。三棵树嵌套嫁接，无论先接哪条内部边，得到的树及每个顶点装饰相同；在两个不同叶边处嫁接也互不影响。这给出两种结合律。无顶点的树不增加运算，给出单位；叶标记的双射与嫁接相容，给出等变条件。

	现在把每个顶点装饰 \(e_v\) 换成 \(\theta(e_v)\)，并在 \(Q\) 中沿树复合。这个求值可按顶点数递归定义：无顶点的树送到 \(\mathbf1\)；其余树从靠近根的顶点看，是该顶点的一个运算，其每条输入边上接着一棵顶点数较少的树。先递归求出这些子树的值，再同时代入根顶点的运算。改变根顶点输入的编号时，\(\theta\) 的等变性和 \(Q\) 的等变条件同时重排子树的值，故结果不变。对子树递归应用同一论证，表明求值不依赖全部临时编号；它也显然对装饰线性并保持树同构，因而定义了 \(\widehat\theta\)。

	若把一棵树嫁接到另一棵树的叶边，递归求值先遇到该叶时，恰好将前一棵树的值代入。嵌套的结合律把这个代入移到根部求值，故求值保持插入复合；它还保持无顶点树和叶的重标号。因此 \(\widehat\theta\) 是 operad 态射，满足交换图
	\[
		\begin{tikzcd}[column sep=large]
			E\arrow[r,"\iota"]\arrow[dr,"\theta"']&\mathcal F(E)\arrow[d,dashed,"\exists!\,\widehat\theta"]\\
			&Q.
		\end{tikzcd}
	\]
	图中的斜箭头是对称序列映射，虚线箭头还须保持全部 operad 结构。每棵装饰树都由单顶点树嫁接得到；任一这样的 operad 态射在单顶点树上已由 \(\theta\) 决定，随后在整棵树上也被复合相容性决定，故唯一。
\end{proof}

\begin{definition}\label{b4:def:operad-ideal}
	设 \(k\) 为域，\(P\) 为 \(k\) 线性 operad。一族右 \(S_n\) 稳定的线性子空间 \(\mathfrak a(n)\subseteq P(n)\) 若满足：只要一个插入因子在 \(\mathfrak a\) 中，插入结果也在 \(\mathfrak a\) 中，即
	\[
		\mathfrak a(n)\circ_i P(m)\subseteq\mathfrak a(n+m-1),
		\qquad P(n)\circ_i\mathfrak a(m)\subseteq\mathfrak a(n+m-1),
	\]
	则称 \(\mathfrak a\) 为一个\term[operad-lixiang]{operad 理想}[operadic ideal]。对一族子空间 \(R(n)\subseteq P(n)\)，含有全部 \(R(n)\) 的 operad 理想之交记为 \((R)\)，称为由这些关系生成的理想。
\end{definition}

所有 \(P(n)\) 自身组成一个包含 \(R\) 的理想，所以定义中的交有意义。重标号与两种插入封闭性均在交下保持，故 \((R)\) 确为含 \(R\) 的最小理想。要求内外两种插入都保持理想，是因为一条关系既可以作为较大表达式的一个子运算，也可以继续代入别的运算。

\begin{proposition}\label{b4:prop:operad-quotient}
	设 \(k\) 为域，\(P\) 为 \(k\) 线性 operad，\(\mathfrak a\) 为 operad 理想。各个商空间 \((P/\mathfrak a)(n)=P(n)/\mathfrak a(n)\) 从 \(P\) 继承 operad 结构。对任意 operad 态射 \(f:P\rightarrow Q\)，若 \(f(\mathfrak a)=0\)，则存在唯一 operad 态射 \(\overline f:P/\mathfrak a\rightarrow Q\) 使 \(f=\overline f\pi\)，其中 \(\pi:P\rightarrow P/\mathfrak a\) 为商映射。特别地，从 \(\mathcal F(E)/(R)\) 到 \(Q\) 的态射，等价于对称序列映射 \(E\rightarrow Q\)，且其自由延拓把 \(R\) 中每个关系送到零。
\end{proposition}
\begin{proof}
	因 \(\mathfrak a(n)\) 对右作用稳定，输入置换在各商空间上有定义。若分别把 \(p,q\) 改为 \(p+a,q+b\)，其中 \(a,b\) 在相应的理想分量中，则插入结果的差为
	\[
		(p+a)\circ_i(q+b)-p\circ_iq
		=a\circ_iq+p\circ_ib+a\circ_ib.
	\]
	三项均属于理想，所以 \([p]\circ_i[q]=[p\circ_iq]\) 良定义。结合、等变和单位等式由原等式取商得到，单位为 \([\mathbf1]\)。

	若 \(f(\mathfrak a)=0\)，各分量的向量空间商泛性质给出唯一线性映射 \(\overline f([p])=f(p)\)。这一公式及 \(f\) 的结构相容性说明 \(\overline f\) 保持置换、插入和单位，因此有
	\[
		\begin{tikzcd}[column sep=large]
			P\arrow[r,"\pi"]\arrow[dr,"f"']&P/\mathfrak a\arrow[d,dashed,"\exists!\,\overline f"]\\
			&Q.
		\end{tikzcd}
	\]
	对于最后的断言，先由\cref{b4:thm:free-operad}把 \(E\rightarrow Q\) 唯一延拓为 \(\widehat\theta:\mathcal F(E)\rightarrow Q\)。其核在右作用与两种插入下封闭，因而是 operad 理想。如果 \(\widehat\theta(R)=0\)，则最小性给出 \((R)\subseteq\ker\widehat\theta\)，故可且只能经商分解。反过来，经此商得到的态射必消去 \(R\)。
\end{proof}

以树的顶点数作为\term[operad-quanzhong]{权重}[weight]，得到分解
\[
	\mathcal F(E)=\bigoplus_{w\geq0}\mathcal F^{(w)}(E),
	\qquad \mathcal F^{(1)}(E)=E.
\]
权重零只有一元恒等运算，插入复合把权重相加。若 \(E(n)=0\) 对所有 \(n\ne2\) 成立，则每个顶点都恰有两条输入边。对含 \(w\geq1\) 个顶点、\(n\) 条叶边的树，进入顶点的边一共是 \(2w\) 条；它们由 \(w-1\) 条内部边和 \(n\) 条叶边组成。因此 \(2w=(w-1)+n\)，即 \(n=w+1\)。无顶点的树满足 \(w=0\)、\(n=1\)，也符合这一公式。尤其二元生成时，权重二关系恰好是三元运算之间的关系。
\symidx{\(\mathcal F^{(w)}(E)\)：自由 operad 中含 \(w\) 个顶点的部分}

\begin{figure}[htbp]
	\centering
	\begin{tikzpicture}[scale=0.9]
		\node (p) at (0,0) {\(p\)};
		\node (q) at (-0.8,1) {\(q\)};
		\node (a) at (-1.4,2) {\(1\)};
		\node (b) at (-0.2,2) {\(2\)};
		\node (c) at (0.8,2) {\(3\)};
		\draw (a)--(q)--(p)--(0,-0.7);
		\draw (b)--(q);
		\draw (c)--(p);
		\node at (0,-1.2) {\(p\circ_1q\)};
		\node (pp) at (5,0) {\(p\)};
		\node (qq) at (5.8,1) {\(q\)};
		\node (aa) at (4.2,2) {\(1\)};
		\node (bb) at (5.2,2) {\(2\)};
		\node (cc) at (6.4,2) {\(3\)};
		\draw (aa)--(pp)--(5,-0.7);
		\draw (bb)--(qq)--(pp);
		\draw (cc)--(qq);
		\node at (5,-1.2) {\(p\circ_2q\)};
	\end{tikzpicture}
	\caption{两个二元运算的插入复合}
	\label{b4:fig:binary-operad-trees}
\end{figure}

在\cref{b4:fig:binary-operad-trees}中，输入从叶流向根，根顶点装饰 \(p\)，另一个顶点装饰 \(q\)。左图表示 \(p(q(v_1,v_2),v_3)\)，右图表示 \(p(v_1,q(v_2,v_3))\)。自由 operad 只记录这两种复合，没有把它们宣布相等。取 \(p=q=\mu\) 后再加入二者之差为零的关系，才得到结合律；这一区别将用于下一节的具体代数结构及本附录最后一节的二次关系。

% !TeX root = ../../Book4.tex
% 纲要标题：### G.2 结合、交换、Lie 与 Poisson 运算
% * 状态：初稿已完成；新增范围尚未正式审核。
% #### G.2.1 Ass 与 Com 的运算和关系
% #### G.2.2 Lie 与 Poisson 的运算和关系
% #### G.2.3 代表性代数与低次核验

\section{结合、交换、Lie 与 Poisson 运算}
\label{b4:sec:G02}

同一个二元符号，施加不同的关系后会产生不同的代数理论。例如 \((ab)c=a(bc)\) 允许去掉乘积的括号，而 \(ab=ba\) 进一步允许忘掉输入次序。Operad 将这两种识别分别记录下来：前者改变复合树的形状，后者还改变对称群的作用。以下都在特征零域 \(k\) 上讨论，输入数至少为一；结合乘法与交换乘法均不预设单位元。Operad 的一元单位仍须作用为恒等映射，这与是否存在乘法单位是两个条件。

\subsection{Ass 与 Com 的运算和关系}
\label{b4:subsec:G02:01}

固定 \(n\ge1\)，考虑每个字母 \(x_1,\ldots,x_n\) 恰好出现一次、没有括号的字。它们有 \(n!\) 个；这些字先是形式符号，不是在某个给定代数中计算后得到的元素。

\begin{definition}\label{b4:def:ass-com-operads}
	设 \(k\) 为特征零域。对每个 \(n\ge1\)，以字母 \(x_1,\ldots,x_n\) 各使用一次的全部 \(n!\) 个字为基，所得向量空间记为 \(\operatorname{Ass}(n)\)。右 \(S_n\) 作用将旧标号 \(i\) 改成 \(\sigma^{-1}(i)\)。复合
	\[
		\gamma(w;w_1,\ldots,w_r)\in\operatorname{Ass}(n_1+\cdots+n_r)
	\]
	对 \(w\in\operatorname{Ass}(r)\) 与 \(w_i\in\operatorname{Ass}(n_i)\) 的基字由代入定义：把第 \(i\) 个输入块依次编号为 \(n_1+\cdots+n_{i-1}+1,\ldots,n_1+\cdots+n_i\)，再将 \(w\) 中的 \(x_i\) 换成相应块上的字 \(w_i\)。按多线性延拓，以字 \(x_1\) 为单位，所得 operad 称为\term[jiehe-operad]{结合 operad}[associative operad]。

	令 \(\operatorname{Com}(n)=k c_n\)，其中 \(S_n\) 作用平凡，并规定
	\[
		\gamma(c_r;c_{n_1},\ldots,c_{n_r})=c_{n_1+\cdots+n_r},
		\qquad \mathbf1=c_1.
	\]
	所得 operad 称为\term[jiaohuan-operad]{交换 operad}[commutative operad]。
\end{definition}
\symidx{\(\operatorname{Ass}\)：结合 operad}
\symidx{\(\operatorname{Com}\)：交换 operad}

例如
\[
	\gamma(x_2x_1;x_2x_1,x_1x_3x_2)=x_3x_5x_4x_2x_1.
\]
这里第一个输入块仍占标号 \(1,2\)，第二个输入块占标号 \(3,4,5\)；外层字 \(x_2x_1\) 使第二个块先出现。字的两次代入，无论先展开内层还是外层，最终都按照同一份嵌套次序排列全部字母，因此满足复合的结合律。代入单字 \(x_1\) 不改变原字。对输入重新编号再代入，与先代入再重新编号所得的字相同，故等变性也成立。\(\operatorname{Com}\) 的这些公理则归结为整数加法及标量乘法的结合律。

\ProseParagraphBreak
若将 \(\sigma\in S_n\) 对应于字
\[
	w_\sigma=x_{\sigma^{-1}(1)}\cdots x_{\sigma^{-1}(n)},
\]
则 \(w_\sigma\cdot\tau=w_{\sigma\tau}\)。因此 \(\operatorname{Ass}(n)\cong k[S_n]\) 是正则右模。这里取逆使记号与\subsecref{b4:subsec:G01:01}的右作用约定一致；字的形式仍更直接地显示运算中的输入次序。

\begin{proposition}\label{b4:prop:ass-com-algebras}
	设 \(k\) 为特征零域，\(A\) 为 \(k\) 向量空间。给定 \(A\) 上的 \(\operatorname{Ass}\) 代数结构，等价于给定一个不预设单位元的双线性结合乘法；给定 \(\operatorname{Com}\) 代数结构，等价于给定一个不预设单位元的双线性交换结合乘法。
\end{proposition}
\begin{proof}
	在 \(\operatorname{Ass}(2)\) 中记 \(\mu=x_1x_2\)。所有字均可由 \(\mu\)、一元单位及输入置换复合得到。两种三元复合满足
	\[
		\mu\circ_1\mu=x_1x_2x_3=\mu\circ_2\mu.
	\]
	所以 operad 态射 \(\theta:\operatorname{Ass}\to\operatorname{End}_A\) 必将 \(\mu\) 送到结合乘法 \(m\)，并把每个字送到按该字次序相乘的多线性映射。反过来，给定结合乘法，定义
	\[
		\theta(x_{i_1}\cdots x_{i_n})(a_1,\ldots,a_n)
		=a_{i_1}\cdots a_{i_n}.
	\]
	结合律使右边与加括号方式无关，也保证字的代入对应映射的复合。换输入次序与字母重编号一致，一元字作用为 \(\operatorname{id}_A\)，因而得到 operad 态射。两个构造互为逆。

	对 \(\operatorname{Com}\)，元素 \(c_2\) 在换位下不变，等变性于是要求 \(m(a,b)=m(b,a)\)；关系 \(c_2\circ_1c_2=c_2\circ_2c_2\) 要求结合律。反过来，交换结合乘法使 \(c_n\) 作用为 \(a_1\cdots a_n\)，并满足全部代入关系。这些构造都没有指定一个 \(0\) 元运算，所以没有产生乘法单位。
\end{proof}

这也给出两种 operad 的生成关系表示。在自由 operad 中取一个无对称限制的二元生成元，即令生成的右 \(S_2\) 模为 \(k[S_2]\)，再模掉 \(\mu\circ_1\mu-\mu\circ_2\mu\)，便得到 \(\operatorname{Ass}\)。确切地说，结合关系能将每棵二叉树改写为按同一次序相乘的右嵌套式；不同输入次序在字模型中线性无关，所以没有遗漏其他关系。

\ProseParagraphBreak
把每个字都送到 \(c_n\)，得到满态射
\[
	q:\operatorname{Ass}\longrightarrow\operatorname{Com}.
\]
其核是由 \(\mu-\mu\cdot(12)\) 生成的 operad 理想。事实上，结合关系允许将一对相邻输入先相乘，交换关系便能交换它们；相邻换位生成全部置换，因此商中所有 \(n\) 元字相等。它们的共同像 \(c_n\ne0\)，所以每个输入数的商恰为一维。\(q\) 所做的是增加交换关系；预复合 \(q\) 则把一个 \(\operatorname{Com}\) 代数看成 \(\operatorname{Ass}\) 代数。

\ProseParagraphBreak
一般自由代数公式\cref{b4:prop:free-operad-algebra}在这里恢复熟悉的对象。对向量空间 \(V\)，有
\[
	\begin{aligned}
		\operatorname{Ass}(V)&\cong\bigoplus_{n\ge1}V^{\otimes n}=T^+(V),\\
		\operatorname{Com}(V)&\cong\bigoplus_{n\ge1}\operatorname{Sym}^n(V)=\operatorname{Sym}^+(V).
	\end{aligned}
\]
第一个同构把 \(w_\sigma\otimes v\) 送到 \(\sigma v\)，利用的是 \(k[S_n]\otimes_{k[S_n]}V^{\otimes n}\cong V^{\otimes n}\)；第二个同构把张量模掉输入置换造成的差，因此得到对称幂。复合在第一式中是张量串接，在第二式中是对称乘法。上标 \(+\) 表示只取正次数部分：没有来自 \(V^{\otimes0}=k\) 的空字。

\subsection{Lie 与 Poisson 的运算和关系}
\label{b4:subsec:G02:02}

结合关系只含一种运算。Poisson 代数则要求一种乘法和一种括号同时存在，还要规定两者怎样相容。自由 operad 的商允许把这三类资料分别写清楚。

\begin{definition}\label{b4:def:lie-poisson-operads}
	设 \(k\) 为特征零域。在只由二元生成元 \(\lambda\) 生成的自由 operad 中，令 \(\lambda\cdot(12)=-\lambda\)，并模掉 Jacobi 元
	\[
		[[x_1,x_2],x_3]+[[x_2,x_3],x_1]+[[x_3,x_1],x_2],
	\]
	所生成的 operad 理想，其中 \([a,b]\) 表示生成元 \(\lambda\) 的代入。所得 operad 记为 \(\operatorname{Lie}\)，称为 \term[Lie-operad]{Lie operad}[Lie operad]。

	再取二元生成元 \(\mu,\lambda\)，其对称性分别为
	\[
		\mu\cdot(12)=\mu,\qquad \lambda\cdot(12)=-\lambda.
	\]
	在它们生成的自由 operad 中，模掉乘法的结合关系、括号的 Jacobi 关系及
	\begin{equation}\label{b4:eq:poisson-operad-leibniz}
		[x_1,x_2x_3]-[x_1,x_2]x_3-x_2[x_1,x_3]
	\end{equation}
	所生成的 operad 理想。这里并置表示 \(\mu\)，方括号表示 \(\lambda\)。所得 operad 记为 \(\operatorname{Pois}\)，称为 \term[Poisson-operad]{Poisson operad}[Poisson operad]。
\end{definition}
\symidx{\(\operatorname{Lie}\)：Lie operad}
\symidx{\(\operatorname{Pois}\)：Poisson operad}

生成元的反对称性已经写在生成右 \(S_2\) 模中；Jacobi 与 Leibniz 关系则位于输入数为三的部分。模掉 operad 理想还会模掉这些关系的全部代入和输入置换。因此商并非只保证三个独立字母满足恒等式，而是保证任意复合运算都保持这些恒等式。

\begin{proposition}\label{b4:prop:lie-poisson-algebras}
	设 \(k\) 为特征零域，\(A\) 为 \(k\) 向量空间。\(A\) 上的 \(\operatorname{Lie}\) 代数结构等价于 Lie 括号；\(A\) 上的 \(\operatorname{Pois}\) 代数结构等价于一个不预设单位元的交换结合乘法，以及对该乘法满足 Leibniz 法则的 Lie 括号。
\end{proposition}
\begin{proof}
	由\cref{b4:thm:free-operad}的泛性质，指定自由 operad 到 \(\operatorname{End}_A\) 的态射，只需指定生成元的像，并保持右 \(S_2\) 作用。\(\lambda\) 的像必须满足 \([a,b]=-[b,a]\)；特征零保证这与交替性等价。再由\cref{b4:prop:operad-quotient}，该态射通过 \(\operatorname{Lie}\) 当且仅当 Jacobi 元的像为零，也就是三个任意输入满足 Jacobi 恒等式。

	对 \(\operatorname{Pois}\)，\(\mu\) 的像同理是交换双线性乘法。三个关系族的像为零，依次就是结合律、Jacobi 恒等式和 Leibniz 法则。一旦这些生成关系为零，它们张成的全部 operad 理想也映为零：对零映射再作任意代入，所得仍为零。因此这些条件既必要又充分。
\end{proof}

\Cref{b4:def:poisson-algebra}对 Poisson 代数另要求乘法含幺；当前 operad 忘掉了这项要求。该处的每个 Poisson 代数都给出这里的 \(\operatorname{Pois}\) 代数。如果乘法原本已有单位，Leibniz 法则自动给出
\[
	\{a,1\}=\{a,1\cdot1\}=2\{a,1\},\qquad \{a,1\}=0;
\]
这说明括号与已有单位相容，却不能从中构造一个原来不存在的单位。

\ProseParagraphBreak
三个输入的关系都是两个二元运算的复合关系。这是它们在下一节成为二次 operad 的原因；此处的“二次”数的是树中的生成顶点，不是输入数。\(\operatorname{Pois}\) 的定义也说明，交换乘法与 Lie 括号单独成立还不够，必须另查\eqref{b4:eq:poisson-operad-leibniz}。

\subsection{代表性代数与低次核验}
\label{b4:subsec:G02:03}

若 \(A\) 为结合代数，则交换子 \([a,b]=ab-ba\) 给出 Lie 括号。反对称性直接成立，而
\[
	\begin{aligned}
		[[a,b],c]&=abc-bac-cab+cba,\\
		[[b,c],a]&=bca-cba-abc+acb,\\
		[[c,a],b]&=cab-acb-bca+bac
	\end{aligned}
\]
的和为零。这些等式在形式字中已经成立，所以有 operad 态射
\[
	\iota:\operatorname{Lie}\longrightarrow\operatorname{Ass},
	\qquad \lambda\longmapsto\mu-\mu\cdot(12).
\]
预复合 \(\iota\) 把结合代数变成其交换子 Lie 代数。箭头从记录 Lie 运算的 operad 指向记录结合运算的 operad，因为我们用结合乘法的表达式实现了 Lie 括号；所得代数结构的限制方向则相反。

\ProseParagraphBreak
任意 Lie 代数配上零乘法都是 \(\operatorname{Pois}\) 代数，因为 Leibniz 法则两边均为零；任意交换结合代数配上零括号也是 \(\operatorname{Pois}\) 代数。两种运算都不为零的例子是 \(k[q,p]\)，其中
\[
	\{f,g\}=\frac{\partial f}{\partial q}\frac{\partial g}{\partial p}
	-\frac{\partial f}{\partial p}\frac{\partial g}{\partial q}.
\]
\Cref{b4:prop:commuting-derivation-star-product}及\eqref{b4:eq:standard-symplectic-poisson}已验证它是 Poisson 括号；这里因此得到一个具体的 \(\operatorname{Pois}\) 代数。对这一已有例子，operad 记录的是普通乘法和括号及其恒等式，不包含星积的全部高阶系数。

\begin{proposition}\label{b4:prop:operad-low-arities}
	设 \(k\) 为特征零域。上述四个 operad 的低次维数为
	\[
		\begin{array}{c|rrrr}
			&\operatorname{Ass}&\operatorname{Com}&\operatorname{Lie}&\operatorname{Pois}\\\hline
			n=2&2&1&1&2\\
			n=3&6&1&2&6
		\end{array}
	\]
	其中 \(\operatorname{Lie}(3)\) 的一组基为
	\[
		\ell_1=[[x_1,x_2],x_3],\qquad
		\ell_2=[[x_2,x_3],x_1].
	\]
	\(\operatorname{Pois}(3)\) 的一组基为
	\[
		x_1x_2x_3,\quad [x_1,x_2]x_3,\quad [x_1,x_3]x_2,
		\quad [x_2,x_3]x_1,\quad \ell_1,\quad\ell_2.
	\]
\end{proposition}
\begin{proof}
	\(\operatorname{Ass}\) 与 \(\operatorname{Com}\) 的维数由定义给出。\(\operatorname{Lie}\) 的二元生成空间是一维，\(\operatorname{Pois}\) 的二元生成空间是对称及反对称的两份一维空间。三元关系及其代入不能降低输入数，所以没有二元生成元因此消失。

	反对称性将任意三元嵌套括号化为
	\[
		\ell_1,\quad \ell_2,\quad \ell_3=[[x_3,x_1],x_2]
	\]
	的线性组合；Jacobi 关系给出 \(\ell_3=-\ell_1-\ell_2\)。为证明剩余两个元素独立，将它们经 \(\iota\) 送入字模型。所得为
	\[
		\begin{aligned}
			\iota(\ell_1)&=x_1x_2x_3-x_2x_1x_3-x_3x_1x_2+x_3x_2x_1,\\
			\iota(\ell_2)&=x_2x_3x_1-x_3x_2x_1-x_1x_2x_3+x_1x_3x_2.
		\end{aligned}
	\]
	在 \(a\iota(\ell_1)+b\iota(\ell_2)\) 中，字 \(x_2x_1x_3\)、\(x_2x_3x_1\) 的系数分别为 \(-a,b\)，所以只能有平凡线性关系。这也证明 \(\iota\) 在输入数为三时为单射。

	对 \(\operatorname{Pois}(3)\)，按复合中括号的个数分成零个、一个、两个三部分。所有定义关系对此次数齐次，故商仍是这三部分的直和。无括号时，结合律与交换律给出单一生成元 \(x_1x_2x_3\)；它在通常的交换代数中可取非零值，因此不为零。有一个括号时，Leibniz 法则把括号在外的式子展开，再用乘法交换性与括号反对称性，得到所列三个“括号乘剩余输入”的生成元。在 \(k[q,p]\) 中分别代入
	\[
		(x_1,x_2,x_3)=(q,p,1),\quad(q,1,p),\quad(1,q,p),
	\]
	三次得到这三个生成元的值依次为 \((1,0,0)\)、\((0,1,0)\)、\((0,0,1)\)。因此它们线性无关。

	两个括号的部分由 \(\ell_1,\ell_2\) 生成。把自由结合代数的交换子 Lie 代数配上零乘法，就得到一个 \(\operatorname{Pois}\) 代数；在其中评价刚才的两个嵌套括号，仍得到两组线性无关的字。因此它们在 \(\operatorname{Pois}(3)\) 中也独立，合起来共有六个基元素。
\end{proof}

这六个三元 Poisson 运算分别对应三种分组：三个输入全部相乘，将其中两个先作括号再乘第三个，以及三个输入全放在嵌套括号中。Leibniz 关系负责把乘法从括号内部展开，Jacobi 关系负责消去嵌套括号之间的冗余。低次计算因而能直接显示各类关系的作用；\(\operatorname{Ass}(3)\) 与 \(\operatorname{Pois}(3)\) 虽然同为六维，它们的复合与对称群作用却并未由这个数字决定。

% !TeX root = ../../Book4.tex
% 纲要标题：### G.3 二次 operad 与 Koszul 对偶
% * 状态：初稿已完成；新增范围尚未正式审核。
% #### G.3.1 二次生成元、关系与对偶
% #### G.3.2 分次、悬移与符号
% #### G.3.3 Ass、Com 与 Lie 的对偶例子
% ---

\section{二次 operad 与 Koszul 对偶}
\label{b4:sec:G03}

结合律连接两次乘法，Jacobi 恒等式连接两次括号，Leibniz 恒等式连接一次乘法和一次括号。它们都含两个运算顶点，因此可以在自由 operad 的同一层中比较。二次对偶把这一层的关系空间换成其正交空间；悬移决定配对所需的符号，使结合关系仍对偶为结合关系。二次关系的计算由有限维线性代数完成，Koszul 性则还要求控制树 bar 复形在全部输入数下的同调。

\subsection{二次生成元、关系与对偶}
\label{b4:subsec:G03:01}

自由 operad 的树形构造有两种不同的计数：叶子数表示输入个数，顶点数表示使用了多少次生成运算。若一个顶点可以有任意多个输入，这两种计数并不相同；若所有生成运算都是二元的，含 \(w\) 个顶点的树恰有 \(w+1\) 个输入。

\begin{definition}\label{b4:def:binary-quadratic-operad}
	设 \(k\) 为特征零的域，\(E\) 为有限维右 \(k[S_2]\) 模，将它看成只在元数二非零的对称序列。记 \(\mathcal F(E)\) 为自由 operad，沿用以顶点数为权重的分次。若 \(R\subseteq\mathcal F^{(2)}(E)(3)\) 是右 \(S_3\) 不变子空间，则
	\[
		\mathcal P=\mathcal F(E)/(R)
	\]
	称为\term[eryuan-erci-operad]{二元二次 operad}[binary quadratic operad]；这里 \((R)\) 是\cref{b4:def:operad-ideal}中由 \(R\) 生成的 operad 理想，商结构由\cref{b4:prop:operad-quotient}给出。
\end{definition}

这一定义中的“二次”指的是运算顶点，不是输入变量的次数。结合律和 Jacobi 恒等式含三个输入，但都是二次关系。关系齐次使商仍按权重分次，并且 \(\mathcal P(1)=k\)、\(\mathcal P(n)=\mathcal P^{(n-1)}(n)\) 对 \(n\ge1\) 成立。若还要给乘法指定零元运算所表示的单位，单位关系会涉及不同权重，不属于这里的二元二次范围。

\ProseParagraphBreak
为了明确关系所在的空间，把两顶点树中的内子树放在根的第一个输入，并按下列三种次序标号：
\[
	(1,2;3),\qquad(2,3;1),\qquad(3,1;2).
\]
例如 \((1,2;3)\) 表示 \(e(f(x_1,x_2),x_3)\)，其中根和内顶点分别装饰 \(e,f\in E\)。这给出向量空间同构
\begin{equation}\label{b4:eq:binary-weight-two}
	\mathcal F^{(2)}(E)(3)\cong(E\otimes E)_{12;3}\oplus(E\otimes E)_{23;1}\oplus(E\otimes E)_{31;2}.
\end{equation}
任意两顶点二叉树的内顶点恰接一对叶子，因此先确定其中一份；若内叶次序相反，就把 \((12)\) 作用于内装饰；若内子树原在根的第二个输入，就把 \((12)\) 作用于根装饰。这样得到唯一的上述代表，说明三份直和中没有遗漏局部输入重排关系。

\begin{definition}\label{b4:def:quadratic-dual-operad}
	设 \(k\) 为特征零的域，\(\mathcal P=\mathcal F(E)/(R)\) 为二元二次 operad，\(E\) 有限维。令
	\[
		E^\vee=E^*\otimes\operatorname{sgn}_2,
	\]
	其中 \(\operatorname{sgn}_2\) 是 \(S_2\) 的一维符号表示，\(E^*\) 取对偶右作用 \((\lambda\cdot\sigma)(e)=\lambda(e\cdot\sigma^{-1})\)。固定符号表示的基 \(\varepsilon\)。在\eqref{b4:eq:binary-weight-two}的三份直和中，规定不同份相互正交，同一份的配对为
	\[
		\left\langle(\lambda\otimes\varepsilon)\otimes(\eta\otimes\varepsilon),\ e\otimes f\right\rangle
		=\lambda(e)\eta(f).
	\]
	记 \(R^\perp\subseteq\mathcal F^{(2)}(E^\vee)(3)\) 为 \(R\) 的零化空间，定义 \(\mathcal P\) 的\term[erci-operad-duiou]{二次对偶 operad}[quadratic dual operad]为
	\[
		\mathcal P^!=\mathcal F(E^\vee)/(R^\perp).
	\]
\end{definition}
\symidx{\(E^\vee=E^*\otimes\operatorname{sgn}_2\)：二元生成空间的带符号对偶}
\symidx{\(\mathcal P^{"!}\)：二元二次 operad 的二次对偶}

上述配对取值于 \(k\)，却不是通常意义下同时重标号不变的配对；它满足
\begin{equation}\label{b4:eq:quadratic-pairing-sign}
	\langle u\cdot\sigma,v\cdot\sigma\rangle=\operatorname{sgn}(\sigma)\langle u,v\rangle
	\qquad(\sigma\in S_3).
\end{equation}
事实上，三循环把三份直和循环交换，不改变内叶方向，因而保持配对；一个换位把内叶方向反转，\(E^\vee\) 中的符号表示恰多给一个负号，\(E^*\) 与 \(E\) 的对偶作用则相消。三循环与换位生成 \(S_3\)，便得到公式。由此 \(R^\perp\) 仍为 \(S_3\) 不变子空间，二次对偶的定义确实给出一个 operad。这个构造与其经典例子见 Ginzburg--Kapranov\cite[Section 2.1]{GinzburgKapranov1994KoszulOperads}。

\begin{proposition}\label{b4:prop:quadratic-operad-bidual}
	设 \(k\) 为特征零的域，\(\mathcal P=\mathcal F(E)/(R)\) 为二元二次 operad，且 \(E\) 有限维。则存在保持生成空间的自然同构 \((\mathcal P^!)^!\cong\mathcal P\)。
\end{proposition}
\begin{proof}
	有限维求值同构给出 \(E^{**}\cong E\)，两次符号表示的张量积为平凡表示，所以 \((E^\vee)^\vee\cong E\)。\eqref{b4:eq:binary-weight-two}中每份配对都是有限维非退化配对，因此第二次取零化空间恢复 \((R^\perp)^\perp=R\)。自由构造和关系商随这两个同构对应，得到所述同构。
\end{proof}

有限维性在这里用于把双对偶认回生成空间，并保证正交空间的双重操作恢复原关系。这一命题只谈生成元与关系，尚未断言任何消解是正合的；后面的 Koszul 性需要另行判断。

\subsection{分次、悬移与符号}
\label{b4:subsec:G03:02}

前面的 operad 集中在次数零。把运算放入分次向量空间时，交换两个齐次因子须使用\cref{b4:exm:graded-braiding}中 \(q=-1\) 的规则：交换次数 \(a,b\) 的两个因子产生 \((-1)^{ab}\)。特别是两个奇次因子的交换会改变符号。

\ProseParagraphBreak
本小节记 \(sV=V[-1]\)，即 \((sV)^j=V^{j-1}\)，把元素次数升高一。它与\cref{b4:def:complex-shift}一致：本卷的 \([1]\) 降低元素次数，所以这里升次的 \(s\) 对应 \([-1]\)。暂不配微分。对齐次多重线性映射 \(f,g\)，分次 End operad 的插入是
\begin{equation}\label{b4:eq:graded-operad-insertion}
	\begin{aligned}
		&(f\circ_i g)(v_1,\ldots,v_{m+n-1})\\
		&\quad=(-1)^{|g|(|v_1|+\cdots+|v_{i-1}|)}
		f(v_1,\ldots,v_{i-1},g(v_i,\ldots,v_{i+n-1}),v_{i+n},\ldots,v_{m+n-1}).
	\end{aligned}
\end{equation}
这个符号来自将映射 \(g\) 跨过它前面的输入，而不是附加一种新的反交换假设。嵌套插入的结合律不变；互不相交的两个插入交换先后次序时，另产生 \((-1)^{|g||h|}\)。两种指数都可按“哪个齐次映射跨过了哪些齐次输入”计算。

\begin{proposition}\label{b4:prop:operadic-suspension-sign}
	设 \(k\) 为特征零的域，\(sk\) 为由次数一的元素 \(e\) 生成的一维分次向量空间。令
	\[
		\mathcal S(n)=\operatorname{Hom}_k((sk)^{\otimes n},sk),\qquad
		u_n(e^{\otimes n})=e.
	\]
	则对每个 \(n\ge1\)，\(\mathcal S(n)\) 集中在次数 \(1-n\)，其右 \(S_n\) 作用为符号表示，且
	\begin{equation}\label{b4:eq:suspension-operad-composition}
		u_m\circ_i u_n=(-1)^{(i-1)(n-1)}u_{m+n-1}.
	\end{equation}
\end{proposition}
\symidx{\(\mathcal S(n)=\operatorname{End}_{sk}(n)\)：悬移产生的一维分次 operad}
\begin{proof}
	输入 \(e^{\otimes n}\) 的次数为 \(n\)，输出次数为一，因此 \(|u_n|=1-n\)。相邻两个输入都是奇次，交换它们乘 \(-1\)，所以一般置换乘其符号。把 \(u_n\) 插入第 \(i\) 个位置时，它跨过 \(i-1\) 个次数一的输入；\eqref{b4:eq:graded-operad-insertion}给出 \((-1)^{(1-n)(i-1)}\)，与所写符号相同。
\end{proof}

对分次 operad \(\mathcal P\)，逐元数取 \(\mathcal S(n)\otimes\mathcal P(n)\)，复合时先按上述交换规则重排张量因子，再分别复合，得到其\term[operad-xuanyi]{operad 悬移}[operadic suspension]。其第 \(n\) 元空间是 \(\mathcal P(n)\) 降低 \(n-1\) 次并扭上 \(\operatorname{sgn}_n\)。二元部分尤其会交换对称与反对称的角色。

\ProseParagraphBreak
在最低的非平凡次数，公式给出
\[
	u_2\circ_1u_2=u_3,\qquad u_2\circ_2u_2=-u_3.
\]
因此两棵二元树在悬移后不能使用相同的配对符号。这个计算正是二次对偶中位置符号的来源。二次对偶 \(\mathcal P^!\) 本身仍按普通次数零的 operad 记写；它在对偶生成元中保留 \(\operatorname{sgn}_2\)，再用关系配对记下树复合的负号，不能把它与逐元数向量空间对偶或仅作一次次数平移混为一谈。

\begin{example}\label{b4:exm:suspended-associativity}
	设 \(A\) 是特征零域上的普通结合代数，乘法记为 \(\mu\)，不要求单位。对 \(sA\) 定义次数 \(-1\) 的运算 \(\widetilde\mu(sa,sb)=s(ab)\)。分次插入给出
	\[
		\begin{aligned}
			(\widetilde\mu\circ_1\widetilde\mu)(sa,sb,sc)&=s((ab)c),\\
			(\widetilde\mu\circ_2\widetilde\mu)(sa,sb,sc)&=-s(a(bc)).
		\end{aligned}
	\]
	故结合律在这一表示中写成 \(\widetilde\mu\circ_1\widetilde\mu+\widetilde\mu\circ_2\widetilde\mu=0\)。若 \(A\) 还交换，则 \(\widetilde\mu\cdot(12)=-\widetilde\mu\)，因为两个悬移后的输入为奇次。这里改变的是运算的次数及其表示，原代数的乘法并没有变成不结合。
\end{example}

\subsection{Ass、Com 与 Lie 的对偶例子}
\label{b4:subsec:G03:03}

先在有限维关系空间里计算二次对偶。令 \(\mu\) 是无预设对称性的二元生成元，则 \(E=k[S_2]\)，基为 \(\mu,\mu\cdot(12)\)。选对偶生成元 \(\nu=\mu^*\otimes\varepsilon\)，就有
\[
	\nu\cdot(12)=-(\mu\cdot(12))^*\otimes\varepsilon.
\]
记 \(L=\mu\circ_1\mu\)、\(R=\mu\circ_2\mu\)，在对偶自由 operad 中同样记 \(L^\vee=\nu\circ_1\nu\)、\(R^\vee=\nu\circ_2\nu\)。三输入自由空间的基为 \(L\cdot\sigma,R\cdot\sigma\)，\(\sigma\in S_3\)。按\eqref{b4:eq:binary-weight-two}把内子树移到根的第一槽后，配对成为
\begin{equation}\label{b4:eq:ass-dual-pairing}
	\begin{aligned}
		\langle L^\vee\cdot\sigma,L\cdot\tau\rangle&=\operatorname{sgn}(\sigma)\delta_{\sigma,\tau},\\
		\langle R^\vee\cdot\sigma,R\cdot\tau\rangle&=-\operatorname{sgn}(\sigma)\delta_{\sigma,\tau},\\
		\langle L^\vee\cdot\sigma,R\cdot\tau\rangle&=\langle R^\vee\cdot\sigma,L\cdot\tau\rangle=0.
	\end{aligned}
\end{equation}
第二行的负号来自根装饰的换位；一般叶标号的符号由\eqref{b4:eq:quadratic-pairing-sign}给出。

\begin{proposition}\label{b4:prop:classical-operad-duals}
	在特征零域 \(k\) 上，有二元二次 operad 的同构
	\[
		\operatorname{Ass}^!\cong\operatorname{Ass},\qquad
		\operatorname{Com}^!\cong\operatorname{Lie},\qquad
		\operatorname{Lie}^!\cong\operatorname{Com}.
	\]
\end{proposition}
\begin{proof}
	结合关系由 \((L-R)\cdot\sigma\)、\(\sigma\in S_3\) 张成。对偶元素写成 \(\sum_\sigma(a_\sigma L^\vee\cdot\sigma+b_\sigma R^\vee\cdot\sigma)\)；根据\eqref{b4:eq:ass-dual-pairing}，它与全部结合关系正交，当且仅当 \(a_\sigma+b_\sigma=0\)。因此对偶关系恰由 \((L^\vee-R^\vee)\cdot\sigma\) 张成，仍是结合律。生成空间 \(E^\vee\) 以 \(\nu\) 及其换位为一组基，仍为正则表示，所以得到 \(\operatorname{Ass}^!\cong\operatorname{Ass}\)。

	对 \(\operatorname{Com}\)，二元生成元对称。在三份循环标号的树上记
	\[
		c_1=(x_1x_2)x_3,\qquad c_2=(x_2x_3)x_1,\qquad c_3=(x_3x_1)x_2.
	\]
	它们在尚未加结合关系的自由 operad 中构成基，关系空间是 \(k(c_1-c_2)+k(c_2-c_3)\)。对偶生成元因符号扭曲而反对称，三份对偶树记为
	\[
		\ell_1=[[x_1,x_2],x_3],\quad \ell_2=[[x_2,x_3],x_1],\quad \ell_3=[[x_3,x_1],x_2].
	\]
	其配对为 \(\langle\ell_i,c_j\rangle=\delta_{ij}\)。正交于两个差的向量必须使三个系数相等，所以对偶关系空间为 \(k(\ell_1+\ell_2+\ell_3)\)，正好是 Jacobi 关系。这证明 \(\operatorname{Com}^!\cong\operatorname{Lie}\)。再用\cref{b4:prop:quadratic-operad-bidual}，或反向取该一维空间的零化空间，得到最后一个同构。
\end{proof}

若把\eqref{b4:eq:ass-dual-pairing}中的负号换成正号，\(L-R\) 的正交空间就由 \(L^\vee+R^\vee\) 张成，得到反结合关系。配对的负号使零化空间仍给出结合关系，因此须与对偶生成元的作用及树的表示一并保留。

\ProseParagraphBreak
二次关系只涉及两顶点树，高输入数下还可能出现“关系之间的关系”。要判断二次资料能否控制这些更高相容条件，需要把所有树组成一个复形。下面给出普通次数零、二元二次情形的 bar 模型。

\begin{definition}\label{b4:def:operadic-bar-complex}
	设 \(k\) 为特征零的域，\(\mathcal P\) 为集中在次数零的二元二次 operad。对 \(n\ge2\)，令 \(B_r(\mathcal P)(n)\) 由下列资料线性生成：有 \(n\) 个标号输入、\(r\) 个顶点的根树，每个顶点至少有两个输入，\(v\) 处装饰 \(\mathcal P(\operatorname{In}(v))\) 中的元素，并给 \(r-1\) 条内部边选择一个顺序 \(e_1,\ldots,e_{r-1}\)。装饰满足多重线性和输入重标号关系；对内部边换序乘该置换的符号。定义
	\begin{equation}\label{b4:eq:operadic-bar-boundary}
		\partial[T;e_1,\ldots,e_{r-1}]
		=\sum_{j=1}^{r-1}(-1)^{j-1}[T/e_j;e_1,\ldots,\widehat e_j,\ldots,e_{r-1}],
	\end{equation}
	其中收缩 \(e_j\) 时，用 operad 的插入复合合并该边两端的装饰。规定 \(B_r(\mathcal P)(n)=0\) 除非 \(1\le r\le n-1\)，并另取 \(B_0(\mathcal P)(1)=k\)，元数一的其他项为零。所得链复形称为此处的\term[operad-bar-fuxing]{operad bar 复形}[operadic bar complex]。
\end{definition}
\symidx{\(B_r(\mathcal P)(n)\)：\(n\) 个输入、链次数 \(r\) 的 operad bar 复形}

这里的 \(r\) 是链次数，\(\partial\) 将其降低一；若改用本卷的上链指标，就令 \(B^q=B_{-q}\)。每个非根顶点唯一对应从它通向根的一条输出内部边，故固定根顶点之后，内部边顺序也给出其余顶点的顺序。这解释了为何定向只列 \(r-1\) 条边，链次数却按全部 \(r\) 个顶点计算；它对应通常 bar 构造中每个顶点各贡献一次悬移。树上装饰原本都是次数零运算，不能把这个链次数与上一小节一个分次运算的内在次数混用。

\begin{proposition}\label{b4:prop:operadic-bar-square-zero}
	对特征零域上集中在次数零的二元二次 operad \(\mathcal P\)，\eqref{b4:eq:operadic-bar-boundary}保持树装饰及边定向关系，且 \(\partial^2=0\)。
\end{proposition}
\begin{proof}
	输入重标号与收缩相容，是 operad 等变公理；装饰的多重线性由复合的双线性保持。边顺序作相邻换位时，未被删除的两边仍贡献一个换位负号；若删去其中一条，则它在交错和中的位置正好相差一，仍改变同一个负号。因此微分在边定向关系下良定义。

	对两条不同内部边 \(e_i,e_j\)、\(i<j\)，先收缩 \(e_i\) 再收缩 \(e_j\) 的符号是 \((-1)^{i-1+j-2}\)，相反次序的符号是 \((-1)^{j-1+i-1}\)，二者相反。两次收缩所得装饰相同：相邻边使用嵌套结合律，不相邻边使用分离输入的结合律。每个二次收缩恰由这两种次序出现，因而逐对相消。
\end{proof}

这个复形在三输入处尤其具体：
\[
	0\longrightarrow\mathcal F^{(2)}(E)(3)
	\xrightarrow{\ \text{复合}\ }\mathcal P(3)\longrightarrow0,
\]
两项分别在链次数二和一。后一项正是前者模掉 \(R\)，所以 \(H_2=R\)、\(H_1=0\)。到四输入时，一棵三顶点树可沿两条内部边作不同收缩；\(\partial^2=0\) 记录了这些关系的相容。更高输入数的正合性不能从三输入计算直接推出。

\begin{definition}\label{b4:def:koszul-operad}
	设 \(k\) 为特征零的域，\(\mathcal P=\mathcal F(E)/(R)\) 为集中在次数零的二元二次 operad，\(E\) 有限维。若对每个 \(n\ge2\) 有
	\[
		H_r(B_\bullet(\mathcal P)(n))=0\qquad(r\ne n-1),
	\]
	则称 \(\mathcal P\) 为\term[Koszul-operad]{Koszul operad}[Koszul operad]。
\end{definition}

这一表述把条件写成 bar 同调只留在最高可能链次数。二元生成保证总权重为 \(n-1\)，故它也是 bar 同调集中在“链次数等于权重”的条件。经典定义及其对偶判据见\cite[Section 4.1]{GinzburgKapranov1994KoszulOperads}。在一般理论中，顶次核带有由切树给出的余复合，二次对偶描述的正是这些资料；Koszul 条件保证可以用这部分较小的资料代替全部 bar 复形。构造二次对偶总是可行，消去其他链次数的同调却是额外结论。

\begin{theorem}\label{b4:thm:classical-koszul-operads}
	设 \(k\) 为特征零的域。则 \(\operatorname{Ass}\)、\(\operatorname{Com}\) 和 \(\operatorname{Lie}\) 都是 Koszul operad。
\end{theorem}
\begin{proofsketch}
	一般判据把 Koszul 性转化为如下条件：对每个自由 \(\mathcal P\) 代数 \(F\)，由二次对偶 \(\mathcal P^!\) 和 \(F\) 上的运算构造的二次标准复形，在把末项 \(F\) 放于零次后，正次数同调全部消失。这里讨论的是由二次资料写出的具体复形。一般理论还说明，在上述有限维二元二次范围内，\(\mathcal P\) 与 \(\mathcal P^!\) 的 Koszul 性等价。判据、对偶性及这三个例子的证明见\cite[Proposition 4.1.4; Theorem 4.2.5; Corollary 4.2.7]{GinzburgKapranov1994KoszulOperads}。

	这里的同调有明确的模型和重编号。结合代数 \(A\) 取约化 bar 复形 \(C_i(A)=A^{\otimes(i+1)}\)、\(i\ge0\)，微分只合并相邻项，不含 Hochschild 链中的首尾环绕项；Lie 代数 \(\mathfrak g\) 取 \(C_i(\mathfrak g)=\Lambda^{i+1}\mathfrak g\)，即删去零次项的 Chevalley--Eilenberg 复形再重编号，使 \(\mathfrak g\) 位于零次。相应同调分别由 \(\operatorname{Tor}_{i+1}^{k\oplus A}(k,k)\) 和 \(\operatorname{Tor}_{i+1}^{U(\mathfrak g)}(k,k)\) 计算。这里 \(k\oplus A\) 是为 \(A\) 添加单位的增广代数，\(U(\mathfrak g)\) 是张量代数模掉 \(xy-yx-[x,y]\) 这些关系所得的包络代数。

	当 \(A=T^+(V)\) 为自由非幺结合代数时，增广代数为 \(T(V)\)。当 \(\mathfrak g\) 为 \(V\) 上的自由 Lie 代数时，其包络代数也为 \(T(V)\)：从两者到结合代数的同态都由从 \(V\) 出发的线性映射唯一确定。两种计算于是都有短自由消解
	\[
		0\longrightarrow T(V)\otimes V\xrightarrow{\ a\otimes v\mapsto av\ }T(V)
		\xrightarrow{\ \epsilon\ }k\longrightarrow0.
	\]
	每个非空张量字都有唯一的末字母，所以左箭头单射、像为增广理想。因此
	\[
		\operatorname{Tor}_j^{T(V)}(k,k)=0\qquad(j\ge2),
	\]
	重编号后的上述同调在正次消失。一般判据给出结合与 Lie 两例，再由 \(\operatorname{Com}^!\cong\operatorname{Lie}\) 及对偶不变性得到交换例。

	这里省略的是一般自由代数同调判据、Lie 代数标准消解的正合性，以及这些消解与树 bar 模型的比较。这些比较是把自由代数的消失计算提升为全部输入数下结论的实质步骤，不能只凭二次对偶的三个同构代替。
\end{proofsketch}

本节的“二次”始终按运算顶点计数。它与二次结合代数按代数生成元次数定义的 Koszul 对偶有相似机制，却属于不同的对象范畴；例如 \(\operatorname{Com}^!=\operatorname{Lie}\) 是两类运算关系之间的对偶，并不是断言某个多项式环的通常二次对偶是 Lie 代数。


\begin{chaptersummary}
	Operad 把一个代数中的运算、输入置换和代入相容性单独抽出。自由树形构造保存全部复合，关系理想保证关系在更大的表达式中仍成立；到自同态 operad 的态射才把抽象运算变成具体多线性映射。由此同一种构造分别给出结合、交换、Lie 和 Poisson 代数。

	\ProseParagraphBreak
	二元二次对偶发生在三输入的关系空间中。生成空间的符号表示扭曲与两种二元插入位置的配对符号，使结合关系仍对偶为结合关系，而交换结合关系与 Jacobi 关系互相对应。树 bar 复形继续比较所有输入数下的关系；Koszul 性要求其同调集中在最高链次数。这比写出二次对偶更强，低次计算能够检验符号和局部相容性，却不能代替对全部输入数的证明。
\end{chaptersummary}

\BookExercises

\begin{exercise}\label{b4:ex:unary-operad-modules}
	设 \(k\) 为域，\(R\) 为含幺结合 \(k\) 代数。令 \(P_R(1)=R\)、\(P_R(n)=0\) 对 \(n\geq2\) 成立，一元复合为 \(r\circ_1s=rs\)，恒等运算为 \(1_R\)。证明这给出一个 operad，并证明 \(P_R\) 代数及其同态恰为左 \(R\) 模及其同态。计算向量空间 \(V\) 上的自由 \(P_R\) 代数。再设对称序列 \(E\) 只在元数一非零，\(E(1)=W\)，证明 \(\mathcal F(E)(1)\) 同构于张量代数 \(T(W)\)，其余元数部分均为零。
\end{exercise}
\begin{solution}
	所有可能非零的插入都在元数一中。两层插入的结合律就是 \((rs)t=r(st)\)，单位律就是 \(1_Rr=r=r1_R\)，而 \(S_1\) 只有恒等置换；因此 \(P_R\) 满足 operad 公理。

	给定 operad 态射 \(P_R\rightarrow\operatorname{End}_V\)，其唯一非零分量是保持乘法和单位的 \(k\) 代数同态 \(\rho:R\rightarrow\operatorname{End}_k(V)\)。规定 \(rv=\rho(r)v\)，得到 \((rs)v=r(sv)\) 和 \(1_Rv=v\)，正是左 \(R\) 模。反过来，左模作用给出这一同态。线性映射 \(h:V\rightarrow V'\) 保持全部 operad 运算，等价于 \(h(rv)=rh(v)\)，即 \(R\) 线性。由自由代数公式，\(P_R(V)=R\otimes_kV\)，作用为 \(s(r\otimes v)=sr\otimes v\)，生成映射为 \(v\mapsto1_R\otimes v\)。

	只含一元顶点的根树只能是一条链，故只有一个叶。含 \(j\) 个顶点的装饰空间为 \(W^{\otimes j}\)，从根向叶排列这些装饰；无顶点的树给出 \(k\)。所以
	\[
		\mathcal F(E)(1)=\bigoplus_{j\geq0}W^{\otimes j}=T(W).
	\]
	把第二条链接到第一条链的输入端，正是把第二组张量接在第一组后面，因而 operad 复合对应张量代数的乘法。没有其他元数的树。这也说明自由 operad 即使生成元空间有限维，一元部分仍可能是无限维的。
\end{solution}

\begin{exercise}\label{b4:ex:transport-operad-algebra}
	设 \(k\) 为域，\(u:V\rightarrow W\) 为向量空间同构。对 \(f\in\operatorname{End}_V(n)\)，令
	\[
		u_*(f)=u\circ f\circ(u^{-1})^{\otimes n}.
	\]
	证明 \(u_*:\operatorname{End}_V\rightarrow\operatorname{End}_W\) 为 operad 同构，并据此把任意 \(P\) 代数结构从 \(V\) 转移到 \(W\)。对 \(u=c\,1_V\)、\(c\in k^\times\)，写出 \(u_*\) 在元数 \(n\) 上的作用。若 \(c=-1\) 且 \(\operatorname{char}k\ne2\)，一个非零结合乘法与它的相反数为何仍定义同构的、不预设单位元的结合代数？
\end{exercise}
\begin{solution}
	对任意 \(w_1,\ldots,w_n\)，\(u_*(f)\) 的值是
	\[
		u\bigl(f(u^{-1}w_1,\ldots,u^{-1}w_n)\bigr).
	\]
	先重排 \(w_i\) 再使用此式，与先重排 \(f\) 的输入再使用此式相同，故 \(u_*\) 保持右 \(S_n\) 作用。计算 \(u_*(f)\circ_i u_*(g)\) 时，内层 \(u\) 与送入外层第 \(i\) 个输入的 \(u^{-1}\) 抵消，得到
	\[
		u_*(f)\circ_i u_*(g)=u_*(f\circ_i g).
	\]
	此外 \(u_*(1_V)=1_W\)。用 \(u^{-1}\) 作同样构造得到逆映射，故这是 operad 同构。若原结构为 \(\rho:P\rightarrow\operatorname{End}_V\)，新结构就是 \(u_*\rho\)，而 \(u\) 本身保持每项运算，成为这两个 \(P\) 代数间的同构。

	对 \(u=c\,1_V\)，多线性给出 \(u_*(f)=c^{1-n}f\)。尤其 \(c=-1\) 时，二元运算变为 \(-f\)。若 \(\mu\) 满足结合律，\(-\mu\) 的两边各含两次乘法，因此仍满足结合律。映射 \(-1_V\) 满足
	\[
		(-1_V)(\mu(a,b))=(-\mu)(-a,-b),
	\]
	故给出所述同构。这里并不要求指定乘法单位；如果原乘法有单位 \(e\)，新乘法的单位相应变为 \(-e\)。
\end{solution}

\begin{exercise}\label{b4:ex:binary-free-bracketings}
	设 \(k\) 为域，对称序列 \(E\) 只在元数二非零，\(E(2)=k[S_2]\) 为正则右模。令 \(P=\mathcal F(E)\)，考虑一个生成向量 \(x\) 上的自由 \(P\) 代数 \(P(kx)\)。证明出现 \(n\) 次 \(x\) 的全部完整加括号的二元词组成一个基。写出 \(n=4\) 时的基；若其基元素个数记为 \(c_n\)，证明
	\[
		c_1=1,\qquad c_n=\sum_{i=1}^{n-1}c_i c_{n-i}\quad(n\geq2),
	\]
	并求 \(c_2,c_3,c_4,c_5\)。
\end{exercise}
\begin{solution}
	正则右 \(S_2\) 模的一个基元素可选定顶点的左、右输入，另一个基元素就是交换这两项后的装饰。把各顶点的装饰吸收到这个左右选择中，装饰树便唯一写成叶有标号的平面二元树。对一个生成向量，\((kx)^{\otimes n}\) 上的 \(S_n\) 作用平凡；自由 \(P\) 代数公式中的平衡关系只把同一平面树的不同叶编号识别。因此其基正是忘记叶编号、每片叶均写 \(x\) 的平面二元树，也就是全部加括号的二元词。

	也可直接检验这里没有隐藏的线性关系：以这些词为形式基构成向量空间，定义两基词的乘积为外加一对括号后的新词，并双线性延拓。这给出任意二元运算的一个代数；把 \(x\) 送到任何带双线性运算的向量空间中的指定元素，逐括号求值即得唯一同态。它因而实现同一个自由对象，各形式基词线性无关。

	四次部分的五个基元素是
	\[
		((xx)x)x,\qquad (x(xx))x,\qquad (xx)(xx),\qquad x((xx)x),\qquad x(x(xx)).
	\]
	对一个含 \(n\geq2\) 片叶的平面二元树，其根的左子树有唯一确定的叶数 \(i\)，右子树有 \(n-i\) 片叶。给定这两棵子树便唯一得到原树，且不同的 \(i\) 不会重复计数，故有题中递推。依次得到
	\[
		c_2=1,\qquad c_3=2,\qquad c_4=5,\qquad c_5=14.
	\]
	它们尚未满足结合律；例如三次部分的 \((xx)x\) 与 \(x(xx)\) 是两个不同的基向量。
\end{solution}

\begin{exercise}\label{b4:ex:free-operad-coproduct}
	设 \(k\) 为域，\(P\) 为只含正元数的 \(k\) 线性 operad。证明零向量空间具有唯一 \(P\) 代数结构，并且是 \(P\) 代数范畴的初对象与终对象。对向量空间 \(V,W\)，证明 \(P(V\oplus W)\) 连同两包含映射诱导的同态，是 \(P(V)\) 与 \(P(W)\) 在 \(P\) 代数范畴中的余积。说明这个结论没有把该余积的底层向量空间断言为 \(P(V)\oplus P(W)\)。
\end{exercise}
\begin{solution}
	零空间上每个正元数的多线性运算都只有一个，且它的一元恒等映射也是零映射，故有唯一 \(P\) 代数结构。任意 \(P\) 代数与它之间的线性映射也唯一。映射到零空间显然保持运算；从零空间映出时，各运算都至少有一个输入，多线性保证其在全零输入处为零，所以也保持运算。这给出初对象与终对象。自由公式中的每项含有 \(0^{\otimes n}=0\)，从而 \(P(0)=0\)。若还编码一个指定的非零常数运算，则这一论证便不适用。

	令 \(A\) 为任意 \(P\) 代数。分别给出同态 \(P(V)\rightarrow A\) 和 \(P(W)\rightarrow A\)，由\cref{b4:prop:free-operad-algebra}，等价于分别给出线性映射 \(V\rightarrow A\) 和 \(W\rightarrow A\)。这又等价于唯一给出线性映射 \(V\oplus W\rightarrow A\)，再由同一泛性质，等价于给出同态 \(P(V\oplus W)\rightarrow A\)。两包含映射的复合恰为原来的两同态，且唯一性也由这一对应得到，故所述对象满足余积的泛性质。

	当 \(P\) 有二元运算时，新对象还允许把来自 \(V\) 和 \(W\) 的生成元混合代入同一运算。例如自由无结合律二元代数中，若 \(v\in V\)、\(w\in W\) 是非零基向量，则形式词 \((vw)\) 既不属于只由 \(V\) 生成的部分，也不属于只由 \(W\) 生成的部分。余积需要这些混合运算，故一般不是两个底层向量空间的直和。
\end{solution}
\begin{exercise}\label{b4:ex:ass-operad-endomorphisms}
	设 \(k\) 为特征零域，\(\mu=x_1x_2\in\operatorname{Ass}(2)\)。求所有 operad 自同态 \(\operatorname{Ass}\to\operatorname{Ass}\)，并判定其中哪些是自同构。可以先考察：给定 \(\alpha,\beta\in k\)，在哪些情形下
	\[
		a*b=\alpha ab+\beta ba
	\]
	对每个结合 \(k\) 代数都定义结合乘法？这里均不要求乘法含幺。
\end{exercise}
\begin{solution}
	Operad 由 \(\mu\) 生成，所以一个自同态完全由
	\[
		\mu\longmapsto\alpha\mu+\beta\mu\cdot(12)
	\]
	决定；它存在当且仅当像满足结合关系。直接展开得
	\[
		(a*b)*c-a*(b*c)
		=\alpha\beta(bac+cab-acb-bca).
	\]
	在自由结合代数中，右边四个字线性无关，因此必须且只需 \(\alpha\beta=0\)。故自同态分成“原乘法的标量倍”和“反向乘法的标量倍”两类，零倍时是同一个自同态。即使二元生成元的像为零，一元单位仍映到一元单位；此时输入数至少为二的部分均映为零。

	当 \(\alpha\ne0,\beta=0\) 时，一个 \(n\) 元字由 \(n-1\) 个二元乘法组成，所以其像为原字的 \(\alpha^{n-1}\) 倍，逆为相应的 \(\alpha^{-1}\) 倍缩放。当 \(\alpha=0,\beta\ne0\) 时，像为反向字的 \(\beta^{n-1}\) 倍；反向两次恢复原字，故也为自同构。除此之外只有二元部分为零的自同态，不可逆。乘法反向与这些缩放交换，因此自同构群同构于 \(k^\times\times C_2\)。
\end{solution}

\begin{exercise}\label{b4:ex:commutative-reflection-ideal}
	设 \(k\) 为特征零域。
	\begin{enumerate}
		\item 对向量空间 \(V\)，证明 \(T^+(V)\) 模掉所有交换子生成的双边理想后，同构于 \(\operatorname{Sym}^+(V)\)。
		\item 在 \(A=M_2(k)\) 中，比较交换子的线性张成空间 \([A,A]\) 与交换子生成的双边理想 \(I\)。求 \(A/[A,A]\) 的向量空间维数及代数 \(A/I\)，并说明为什么 \(A/[A,A]\) 不能在这里代替交换商代数。
	\end{enumerate}
\end{exercise}
\begin{solution}
	1. 对称乘法给出满同态 \(T^+(V)\to\operatorname{Sym}^+(V)\)，交换子理想包含于其核。在商中可交换相邻的两个张量因子；相邻换位生成全部置换，故每个 \(n\) 次齐次部分恰将 \(V^{\otimes n}\) 模掉 \(v-\sigma v\) 这些差。另一方面，这些差均由相邻交换子乘上左右张量得到；左右某一边为空时，只需省去该边乘法。因此商的各次部分正是 \(\operatorname{Sym}^n(V)\)，乘法也是对称乘法。

	2. 矩阵迹零化交换子，故 \([A,A]\) 包含于无迹矩阵空间。矩阵单位满足
	\[
		[E_{11},E_{12}]=E_{12},\qquad
		[E_{22},E_{21}]=E_{21},\qquad
		[E_{12},E_{21}]=E_{11}-E_{22}.
	\]
	右边构成无迹矩阵空间的一组基，故 \([A,A]\) 恰为该三维空间，\(A/[A,A]\) 是一维向量空间。但 \(h=E_{11}-E_{22}\in I\) 且 \(h^2=1\)，所以 \(1\in I\)，进而 \(I=A\) 及 \(A/I=0\)。空间 \([A,A]\) 不是双边理想：它含 \(h\)，却不含 \(h^2=1\)，因为 \(\operatorname{Tr}(1)=2\ne0\)。因此矩阵乘法不能下降到 \(A/[A,A]\)。增加交换关系需要模掉对乘法封闭的理想，不能只模掉交换子的线性张成空间。
\end{solution}

\begin{exercise}\label{b4:ex:poisson-plane-polynomial}
	设 \(k\) 为特征零域，\(A=k[x,y]\)。证明：对任意 \(h\in A\)，公式
	\[
		\{f,g\}_h=h\cdot\left(
		\frac{\partial f}{\partial x}\frac{\partial g}{\partial y}
		-\frac{\partial f}{\partial y}\frac{\partial g}{\partial x}
		\right)
	\]
	给出 Poisson 括号；反过来，\(A\) 上的每个 Poisson 括号都唯一以这种方式出现。证明应说明 \(h\) 不必是常数，以及 Jacobi 恒等式为何仍成立。
\end{exercise}
\begin{solution}
	任何交替双导子都由 \(\{x,x\},\{x,y\},\{y,x\},\{y,y\}\) 决定。具体地，反复使用乘积法则，得到
	\[
		\{x^i y^j,x^r y^s\}
		=(is-jr)\,h x^{i+r-1}y^{j+s-1},
	\]
	其中系数为零的项直接取零，不使用负次幂。由线性性便得到题中公式，且 \(h=\{x,y\}\) 唯一。反过来，所列公式是交替双导子，尚需核验 Jacobi 恒等式。

	对交换代数上的任意交替双导子，记其 Jacobi 表达式为
	\[
		J(a,b,c)=\{a,\{b,c\}\}+\{b,\{c,a\}\}+\{c,\{a,b\}\}.
	\]
	它是三线性的；当两个输入相等时，例如 \(J(a,a,c)\)，前两项因 \(\{a,c\}=-\{c,a\}\) 抵消，第三项为零，所以它还是交替的。它在每个变量上也满足乘积法则。以第一变量为例，展开 \(J(ab,c,d)\)，其中含两个嵌套括号的项合成 \(a\cdot J(b,c,d)+b\cdot J(a,c,d)\)，额外项为
	\[
		\{c,a\}\{d,b\}+\{c,b\}\{d,a\}
		+\{d,a\}\{b,c\}+\{d,b\}\{a,c\}=0.
	\]
	最后一个等号使用反对称性和乘法交换性。循环置换给出另两个变量的乘积法则。

	因此，在 \(k[x,y]\) 上检验 \(J=0\)，只需检验三个输入各取 \(x\) 或 \(y\) 的情形：乘积法则会把任意多项式输入展开为这些值的线性组合。三个输入只有两种选择，必有重复，而交替性使该值为零。这个证明适用于任意 \(h\)，并未要求它的偏导数为零。
\end{solution}

\begin{exercise}\label{b4:ex:poisson-tensor-product}
	设 \(k\) 为特征零域，\(A,B\) 为交换含幺 Poisson \(k\) 代数。证明通常的交换代数 \(A\otimes_kB\) 在下列括号下也是 Poisson 代数：
	\[
		\{a\otimes b,a'\otimes b'\}
		=\{a,a'\}\otimes bb'+aa'\otimes\{b,b'\}.
	\]
	并证明这是唯一使两个映射 \(a\mapsto a\otimes1\)、\(b\mapsto1\otimes b\) 保持括号，且使两份像之间的括号恒为零的 Poisson 括号。由此说明 \(k[q_1,p_1,q_2,p_2]\) 上两个独立标准辛括号怎样合成。
\end{exercise}
\begin{solution}
	右边关于四个变量均为 \(k\) 线性，所以定义了张量积上的双线性映射。反对称性由两因子的反对称性给出。若固定 \(a\otimes b\)，则映射 \(a'\otimes b'\mapsto\{a,a'\}\otimes bb'\) 是“\(A\) 上的导子乘以中心元素 \(1\otimes b\)”，另一项同理；两项均为 \(A\otimes B\) 上的导子，故 Leibniz 法则成立。

	对交替双导子，Jacobi 表达式是每个变量上的导子。其展开理由是：在一个变量中代入乘积后，额外出现的交叉项两两为
	\[
		\{c,a\}\{d,b\}+\{d,b\}\{a,c\}=0,
		\qquad
		\{c,b\}\{d,a\}+\{d,a\}\{b,c\}=0.
	\]
	现在 \(A\otimes B\) 由 \(a\otimes1\) 与 \(1\otimes b\) 这些元素生成。Jacobi 表达式在三个生成元均来自同一因子时由该因子的 Jacobi 恒等式为零；当两个来自一份、另一个来自另一份时，跨两份的括号为零，三个嵌套项也都为零。因此利用三个变量的乘积法则，Jacobi 表达式处处为零。

	唯一性来自 \(a\otimes b=(a\otimes1)(1\otimes b)\)：对两个这样的乘积使用两次 Leibniz 法则，跨因子的两项为零，剩下的两项恰为题中公式。对 \(A=k[q_1,p_1]\)、\(B=k[q_2,p_2]\)，于是得到
	\[
		\{q_i,p_j\}=\delta_{ij},\qquad
		\{q_i,q_j\}=\{p_i,p_j\}=0.
	\]
	多项式的括号由这些生成元关系及乘积法则唯一决定，正是两对变量的标准辛括号。
\end{solution}

\begin{exercise}\label{b4:ex:poisson-quotient-ideals}
	设 \(k\) 为特征零域，\(A\) 为交换含幺 Poisson \(k\) 代数，\(I\) 为其乘法理想。
	\begin{enumerate}
		\item 证明括号由 \(A\) 下降到 \(A/I\)，当且仅当 \(\{A,I\}\subseteq I\)。
		\item 对标准辛 Poisson 代数 \(k[q,p]\)，证明这样的理想只有 \(0\) 和全环。
		\item 在 \(k[x,y]\) 上改取 \(\{f,g\}=x(f_xg_y-f_yg_x)\)。证明 \((x)\) 是满足条件的非零真理想，并求商代数的括号。
	\end{enumerate}
\end{exercise}
\begin{solution}
	1. 如果括号能够下降，\(i\in I\) 的商类为零，故 \(\{a,i\}\) 的商类也必须为零。反之，若 \(\{A,I\}\subseteq I\)，则用 \(a+i,b+j\) 代替 \(a,b\) 后，括号的差
	\[
		\{a,j\}+\{i,b\}+\{i,j\}
	\]
	属于 \(I\)，因此商上的括号良定义。双线性、Jacobi 及 Leibniz 恒等式都由原式取商得到。

	2. 在标准辛括号下，\(\{q,f\}=\partial f/\partial p\)，\(\{p,f\}=-\partial f/\partial q\)。所以满足条件的理想对两个偏导均封闭。若 \(I\ne0\)，在其中取总次数最小的非零多项式 \(f\)。若 \(f\) 非常数，特征零保证它至少有一个偏导非零，该偏导仍在 \(I\) 中且次数更低，矛盾。因此 \(I\) 含非零常数，必为全环。

	3. 此括号的 Jacobi 恒等式由\cref{b4:ex:poisson-plane-polynomial}中 \(h=x\) 的情形成立：Jacobi 表达式对每个变量满足乘积法则，故可将三个输入约化到 \(x,y\)，因必有重复而为零。对任意 \(f,g\)，括号 \(\{f,xg\}\) 的公式整体含因子 \(x\)，故属于 \((x)\)。商代数为 \(k[y]\)，所有括号均为零，因为原括号的每个值都在 \((x)\) 中。这与标准辛括号下没有非零真理想的结论不同。
\end{solution}
\begin{exercise}\label{b4:ex:G-magmatic-dual}
	设 \(k\) 为特征零的域，\(E=k[S_2]\)，令 \(\mathcal M=\mathcal F(E)\)，即不对二元运算施加任何关系。求 \(\mathcal M^!\) 的各元数空间，并说明它的代数满足什么条件。构造一个二维的 \(\mathcal M^!\) 代数，其乘法非零；说明它不可能有乘法单位。
\end{exercise}
\begin{solution}
	因为原关系空间为零，对偶关系空间为整个 \(\mathcal F(E^\vee)^{(2)}(3)\)。因此 \(\mathcal M^!(1)=k\)、\(\mathcal M^!(2)=E^\vee\cong k[S_2]\)，所有两次嵌套乘法均为零。任意至少含两个顶点的二叉树都含一对相邻运算顶点；把该两顶点部分看成一个三元关系，再与剩余部分复合，整树便属于关系理想。所以 \(\mathcal M^!(n)=0\) 对所有 \(n\ge3\) 成立。

	其代数是带任意双线性乘法且满足 \((ab)c=a(bc)=0\) 的向量空间。取基 \(u,v\)，规定 \(u^2=v\)，其余基向量乘积均为零，即得所需例子。若存在单位 \(1\)，则 \((1\cdot1)a=0\) 与单位律给出 \(a=0\) 对每个元素成立；故非零的这种代数不可能有单位。这也说明不含零元运算的约定在对偶构造中有实质作用。
\end{solution}

\begin{exercise}\label{b4:ex:G-scaled-associativity-dual}
	设 \(k\) 为特征零的域，\(q\in k^\times\)。令 \(\mathcal P_q\) 为由一个无预设对称性的二元运算生成、以 \((ab)c=q\,a(bc)\) 及其输入重标号为关系的 operad。证明 \(\mathcal P_q^!\cong\mathcal P_{q^{-1}}\)。若 \(q\ne1\)，证明 \(\mathcal P_q\) 代数中任意四个元素的乘积，不论怎样加括号，都为零；说明这一结论怎样从同一四输入式的两种化简次序产生。
\end{exercise}
\begin{solution}
	用\eqref{b4:eq:ass-dual-pairing}的记号，原关系由 \((L-qR)\cdot\sigma\) 张成。对偶系数 \(a_\sigma L^\vee\cdot\sigma+b_\sigma R^\vee\cdot\sigma\) 正交于该关系，等价于 \(a_\sigma+q b_\sigma=0\)。因此对偶关系由 \((L^\vee-q^{-1}R^\vee)\cdot\sigma\) 张成，得到所述同构。

	同一个四输入式可计算为
	\[
		\begin{aligned}
			((ab)c)d&=q(ab)(cd)=q^2a(b(cd)),\\
			((ab)c)d&=q(a(bc))d=q^2a((bc)d)=q^3a(b(cd)).
		\end{aligned}
	\]
	所以 \(q^2(1-q)a(b(cd))=0\)。当 \(q\ne0,1\) 时系数可逆，右结合的四重乘积为零。五种括号式分别通过所给关系化为右结合式的标量倍，故全部为零。这个消失是二次关系在四输入中的后果，不是另加的四元定义；它也说明写出三输入关系之后，仍须检查高输入数中的相容性。
\end{solution}

\begin{exercise}\label{b4:ex:G-general-shift-operad}
	设 \(k\) 为特征零的域，\(L_d\) 为由次数 \(d\in\mathbb Z\) 的元素 \(e\) 生成的一维分次向量空间。令 \(u_n(e^{\otimes n})=e\)。求 \(\operatorname{End}_{L_d}(n)\) 的次数、\(S_n\) 作用和部分复合，并核验 \((u_2\circ_1u_2)\circ_3u_2\) 与 \((u_2\circ_2u_2)\circ_1u_2\) 之间的分次交换符号。
\end{exercise}
\begin{solution}
	映射 \(u_n\) 的次数为 \(d(1-n)\)。交换两个输入贡献 \((-1)^{d^2}=(-1)^d\)，故右作用为 \(u_n\cdot\sigma=(\operatorname{sgn}\sigma)^d u_n\)。插入 \(u_n\) 时，它跨过 \(i-1\) 个次数为 \(d\) 的输入，因此
	\[
		u_m\circ_i u_n=(-1)^{d^2(i-1)(n-1)}u_{m+n-1}
		=(-1)^{d(i-1)(n-1)}u_{m+n-1}.
	\]
	题中第一个式子为 \(u_3\circ_3u_2=u_4\)，第二个为 \((-1)^d u_3\circ_1u_2=(-1)^d u_4\)。它们是把两个次数 \(-d\) 的二元运算插入根的不同输入时交换先后次序，预期符号为 \((-1)^{(-d)(-d)}=(-1)^d\)，与计算一致。特别地，偶数悬移不改变这些交换符号，奇数悬移则不能省略它们。
\end{solution}

\begin{exercise}\label{b4:ex:G-ass-bar-four-inputs}
	在特征零域上，考虑 \(B_\bullet(\operatorname{Ass})(4)\) 中保持叶子从左至右为 \(1,2,3,4\) 的部分。用 \(a,b,c,d\) 表示这四个输入。三顶点树取
	\[
		\begin{gathered}
			T_1=((ab)c)d,\quad T_2=(a(bc))d,\quad T_3=(ab)(cd),\\
			T_4=a((bc)d),\quad T_5=a(b(cd)).
		\end{gathered}
	\]
	把相应两条内部边按“先从根出发访问子树、同层从左至右”的次序定向。证明这一部分的 bar 复形维数为 \(5\to5\to1\)，写出两张边界矩阵，求出全部同调，并给出顶次同调的一个非零代表。
\end{exercise}
\begin{solution}
	链次数二的五棵树取为 \(A,B,C,D,E\)：\(A\) 是根的第一输入接三元运算，\(B\) 是根的第二输入接三元运算；\(C,D,E\) 是一个三元根分别在第一、第二、第三输入接二元运算。每棵只有一条内部边。链次数一只有四元乘法 \(F\)。所有装饰都保持原输入顺序，故 \(\partial A=\partial B=\partial C=\partial D=\partial E=F\)。

	依照题定边顺序，\eqref{b4:eq:operadic-bar-boundary}给出
	\[
		\begin{aligned}
			\partial T_1&=C-A,&\partial T_2&=D-A,&\partial T_3&=E-C,\\
			\partial T_4&=D-B,&\partial T_5&=E-B.
		\end{aligned}
	\]
	例如 \(T_1\) 先收缩根与其左子树根之间的边，得到三元根第一输入中的二元乘法 \(C\)；再收缩更深的边，得到 \(A\)，交错符号为负。以所列基为列、像空间基为行，两张矩阵为
	\[
		[\partial_3]=
		\begin{pmatrix}
			-1&-1&0&0&0\\
			0&0&0&-1&-1\\
			1&0&-1&0&0\\
			0&1&0&1&0\\
			0&0&1&0&1
		\end{pmatrix},\qquad
		[\partial_2]=\begin{pmatrix}1&1&1&1&1\end{pmatrix}.
	\]
	每列元素之和为零，所以 \(\partial_2\partial_3=0\)。解 \(\partial_3(x_1,\ldots,x_5)^{\mathsf T}=0\) 得
	\[
		(x_1,x_2,x_3,x_4,x_5)=t(1,-1,1,1,-1).
	\]
	故 \(\partial_3\) 的秩为四；\(\ker\partial_2\) 也是四维，且含 \(\operatorname{im}\partial_3\)，二者相等。于是 \(H_1=H_2=0\)，\(H_3\) 一维，由 \(T_1-T_2+T_3+T_4-T_5\) 代表。

	完整的四输入复形还允许全部 \(4!\) 种输入顺序，且边界不改变这个顺序，故分成上述复形的 \(24\) 份直和。其顶次同调维数为 \(24\)，其余次数为零。这是四输入处的直接验证，不代替任意输入数下的 Koszul 性定理。
\end{solution}
